% D28
% Auto-generated from markdown source
% origin: ai

\documentclass[12pt, a4paper]{article}
\usepackage[utf8]{inputenc}
\usepackage[T1]{fontenc}
\usepackage{amsmath, amssymb, amsfonts, amsthm}
\usepackage{textcomp}
\usepackage{geometry}
\usepackage{hyperref}
\usepackage{booktabs}
\usepackage{graphicx}
\usepackage{fancyhdr}
\usepackage{tcolorbox}
\usepackage{enumitem}

\geometry{margin=1in}
\pagestyle{fancy}
\fancyhf{}
\rhead{D28}
\lhead{\thepage}

\newtheorem{definition}{Definition}[section]
\newtheorem{theorem}{Theorem}[section]
\newtheorem{corollary}{Corollary}[section]
\newtheorem{proposition}{Proposition}[section]
\newtheorem{lemma}{Lemma}[section]
\newtheorem{conjecture}{Conjecture}[section]

\newcommand{\calI}{\mathcal{I}}
\newcommand{\calP}{\mathcal{P}}
\newcommand{\calE}{\mathcal{E}}
\newcommand{\calS}{\mathcal{S}}
\newcommand{\calH}{\mathcal{H}}
\newcommand{\calO}{\mathcal{O}}
\newcommand{\calF}{\mathcal{F}}
\newcommand{\calU}{\mathcal{U}}
\newcommand{\calV}{\mathcal{V}}
\newcommand{\calA}{\mathcal{A}}
\newcommand{\calG}{\mathcal{G}}
\newcommand{\calB}{\mathcal{B}}
\newcommand{\calC}{\mathcal{C}}
\newcommand{\calL}{\mathcal{L}}
\newcommand{\calM}{\mathcal{M}}
\newcommand{\calX}{\mathcal{X}}
\newcommand{\calY}{\mathcal{Y}}
\newcommand{\calW}{\mathcal{W}}
\newcommand{\calT}{\mathcal{T}}
\newcommand{\calZ}{\mathcal{Z}}
\newcommand{\calR}{\mathcal{R}}
\newcommand{\calQ}{\mathcal{Q}}
\newcommand{\calD}{\mathcal{D}}
\newcommand{\calN}{\mathcal{N}}
\newcommand{\calK}{\mathcal{K}}

\newcommand{\Tr}{\text{Tr}}
\newcommand{\E}{\mathbb{E}}
\newcommand{\diff}{\mathrm{d}}
\newcommand{\KL}{D_{\mathrm{KL}}}
\newcommand{\argmax}{\arg\max}
\newcommand{\argmin}{\arg\min}
\newcommand{\hbarF}{\hbar}

\newtcolorbox{thmbox}[1]{
  colback=blue!5!white, colframe=blue!75!black, fonttitle=\bfseries,
  title={#1}, breakable
}
\newtcolorbox{defbox}[1]{
  colback=green!5!white, colframe=green!75!black, fonttitle=\bfseries,
  title={#1}, breakable
}

\begin{document}

\title{D28}
\author{RALPH Loop}
\date{}
\maketitle

\begin{abstract}
This paper presents D28, formalizing the underlying principles using information-theoretic methods.

\textbf{Keywords}: D28, information theory, formalization.
\end{abstract}

\tableofcontents
\section{D28}




\subsection{0.}


\textbf{**  }**

\textbf{D28 } C  C   ****



\subsection{1.}


\subsubsection{1.1}


\textbf{ D28-1}
\(\mathcal{S}\) \(D \gg 1\)
**** \(\Pi_i: \mathcal{S} \to \mathbb{R}^{d_i}\)  \(i\) \(d_i \ll D\)

****
\[C_1 \perp C_2 \text{ ( } \Pi_1 \text{ )} \iff \langle \Pi_1 C_1, \Pi_1 C_2 \rangle = 0\]

****
\[\exists \alpha, \beta: \langle \alpha C_1, \beta C_2 \rangle_{\mathcal{S}} \neq 0\]


\subsubsection{1.2}


\textbf{ D28-2}
\(\mathcal{D}_1, \mathcal{D}_2 \subset \mathcal{S}\) **** \(\mathcal{A}: \mathcal{D}_1 \to \mathcal{D}_2\)
\[\|\mathcal{A}(x) - \mathcal{A}(y)\|_{\mathcal{S}} \approx \|x - y\|_{\mathcal{D}_1} \quad \forall x, y \in \mathcal{D}_1\]

****
\[F(\mathcal{A}) = 1 - \frac{1}{|\mathcal{D}_1|^2} \sum_{x,y} \left(\frac{\|\mathcal{A}(x) - \mathcal{A}(y)\|_{\mathcal{S}} - \|x - y\|_{\mathcal{D}_1}}{\|x - y\|_{\mathcal{D}_1}}\right)^2\]


\subsubsection{1.3}


\textbf{ D28-3}
**** \(\xi(t)\)
\[\frac{d}{dt} \theta(t) = -\nabla \Phi(\theta(t)) + \xi(t)\]

**** \(\Delta_{12}\)
\[\Delta_{12}(t) = \mathbb{E}[\langle \Pi_1 \theta(t), \Pi_2 \theta(t) \rangle]\]

****
\[\Delta_{12}(t) > \Delta_{thresh} \land \frac{d\Delta_{12}}{dt} > 0\]


\subsection{2.}


\subsubsection{ D28-1}


****
****
\[\mathcal{P}_{completeness} = \frac{\sum_i \|\Pi_i^* \Pi_i\|}{\|I_{\mathcal{S}}\|}\]

\textbf{D28-1 }
\[P(C_1 \perp C_2 | \text{}) = \frac{\sum_{k=d}^{D} \lambda_k}{\sum_{k=1}^{D} \lambda_k}\]

\(\lambda_k\) \(d\) \(D\)

\textbf{ D28-1a}
\[\lim_{d \to D} P(C_1 \perp C_2 | \text{}) = P(C_1 \perp C_2 | \text{})\]
\[\lim_{d \to 0} P(C_1 \perp C_2 | \text{}) = 1\]

\textbf{**  }**


\subsubsection{ D28-2}


****
**** \(\mathcal{G}_{analog} = (V, E)\) \(V\) \((u,v) \in E\)

\textbf{D28-2 }
\[\mathcal{G}_{analog} \text{  } \mathcal{G}_{semantic} \iff \forall (u,v) \in E_{semantic}: P(\mathcal{A}_{u \to v}) > \epsilon\]

****
\[\eta_{explore} = \frac{\sum_{(u,v) \in E_{analog}} I(u;v)}{\sum_{(u,v) \in V^2} H(u)}\]

\textbf{ D28-2a}
\begin{itemize}
\item ****\(\mathcal{G}_{analog}\)  \(\mathcal{G}_{semantic}\)
\item ****
\end{itemize}



\subsubsection{ D28-3}


****
****
\[d_{analog}(x, y) = \inf_{\mathcal{A} \in \mathcal{A}_{domain}} \|\mathcal{A}(x) - y\|_{\mathcal{S}}\]

\textbf{D28-3 }
\[d_{analog}(x, y) \leq d_{semantic}(x, y)\]
\(x, y\)

****
\[d_{analog}(x, z) \leq d_{analog}(x, y) + d_{analog}(y, z)\]
****


\subsubsection{ D28-4}


****
****
\[E_{cognitive}(\gamma) = \int_{\gamma} \sqrt{g_{ij} \dot{x}^i \dot{x}^j} dt + \lambda \cdot \mathbb{I}(\gamma \text{ crosses domain})\]

\(g_{ij}\)  Fisher ****

\textbf{D28-4 }
\[\arg\min_{\gamma: x \to y} E_{cognitive}(\gamma) = \text{Geodesic}_{analog}(x, y)\]

****
\[\min_{\gamma} E_{cognitive}(\gamma) = \min_{\gamma_1} E_1 + \min_{\gamma_2} E_2\]
\(\gamma = \gamma_1 \cup \gamma_2\) ****

\textbf{ D28-4a}
****


\subsubsection{ D28-5}


****
****
\[I(C_1; C_2) = \mathbb{E}\left[\log \frac{P(C_1, C_2)}{P(C_1)P(C_2)}\right]\]

\textbf{D28-5 }-
\[P(C_1 \perp C_2 | \text{}) = 1 \implies I(C_1; C_2) \leq I_{threshold}(\Pi)\]

****
\[I_{threshold}(\Pi) = H(C_1) \cdot H(C_2) \cdot \left(1 - \frac{\text{rank}(\Pi)}{D}\right)\]

****
\(D \gg \text{rank}(\Pi)\) \(C_1 \perp C_2\)

\textbf{ D28-5a}
****   \(D/d\)


\subsection{3.}


\subsubsection{ D28-A}


\textbf{ D28-1}
1.  word2vec 300D
2.  100 vs
3.  SVD  \(d = 5, 10, 50, 100, 300\)
4.  \(P(C_1 \perp C_2 | d)\)
5. **** 100\%


\subsubsection{ D28-B}


\textbf{ D28-4}
1.  \(\mathcal{G}_{semantic}\)
2.  vs
3. **** \(<\)


\subsubsection{ D28-C}


\textbf{ D28-2}
1.  =
2.
3. ****


\subsubsection{ D28-D-}


\textbf{ D28-5}
1.
2.
3. **** \(I_{threshold} \propto (1 - d/D)\)


\subsection{4.}


\subsubsection{4.1 D24}


\textbf{D28-D24 }
\begin{itemize}
\item D28 \textbf{**  D24 }**
\item D28-3 \textbf{** = D24 }**
\item ****
\end{itemize}



\subsubsection{4.2 D25}


\textbf{D28-D25 }
\begin{itemize}
\item D28 \textbf{**  D25 }**
\item D28-4 \textbf{** = D25-3 }**
\item ****
\end{itemize}



\subsubsection{4.3 D1-S12}


\textbf{D28-D1-S12 }
\begin{itemize}
\item D28 \textbf{** = D1-S12 }**
\item D28-5 \textbf{-}  D1-S12 \textbf{ \(\rho\)}
\item ****
\end{itemize}



\subsection{5.}


\subsubsection{5.1}


1. \textbf{D28-1 }
2. \textbf{D28-4 }
3. \textbf{D28-5 }

\subsubsection{5.2}


1. ****
2. ****
3. ****


\subsection{6.}


\#\#\#

**** \(\mathcal{S}\)  \(C_1, C_2\)  \(\Pi\)

****
1. ****\(P(C_1 \perp C_2 | \text{}) = \sum_{k=d}^{D} \lambda_k / \sum_{k=1}^{D} \lambda_k\)
2. ****\(\mathcal{G}_{analog}\)  \(\mathcal{G}_{semantic}\)
3. ****\(d_{analog}(x,y) \leq d_{semantic}(x,y)\)
4. **** =
5. \textbf{-}\(I(C_1; C_2) \leq H(C_1)H(C_2)(1 - d/D)\)

\textbf{** }\textbf{*}*

\textbf{**  }**


*D28 *

\begin{thebibliography}{99}
\bibitem{ref1} Shannon, C.E. (1948). A Mathematical Theory of Communication.
\bibitem{ref2} Cover, T.M. (2006). Elements of Information Theory.
\end{thebibliography}

\end{document}